\chapter{Medicines and alcohol}

\section*{Definition and Overview}
Medicines and alcohol interactions refer to the complex pharmacokinetic and pharmacodynamic changes that occur when ethanol is consumed concurrently with therapeutic substances. These interactions can significantly alter the therapeutic efficacy, safety profile, and toxicity of both prescribed and over-the-counter medications. Ethanol is a central nervous system (CNS) depressant that is rapidly absorbed from the gastrointestinal tract and metabolised primarily in the liver by alcohol dehydrogenase (ADH) and the cytochrome P450 enzyme system, specifically the CYP2E1 pathway. When combined with medications, alcohol can produce additive CNS depression, modify hepatic drug metabolism through enzyme inhibition (acute consumption) or induction (chronic consumption), and exacerbate gastrointestinal, cardiovascular, and metabolic adverse effects. Understanding these interactions is a critical component of clinical pharmacology and patient safety, as they represent a major cause of preventable drug-related hospitalisations and adverse clinical outcomes.

\section*{Epidemiology}
Concomitant use of alcohol and medicines is highly prevalent globally and within Australia, posing a significant public health concern. Epidemiological surveys indicate that approximately 20\% to 30\% of individuals who drink alcohol also use medications that are known to interact adversely with ethanol. This risk is particularly pronounced in the ageing population; older adults frequently take multiple prescription medications (polypharmacy) and have age-related changes in body composition, renal clearance, and hepatic metabolism, which increases their vulnerability to toxic interactions. Furthermore, registry data show that adverse drug-alcohol interactions are common contributors to emergency department presentations for accidental drug poisonings, falls, motor vehicle accidents, and acute gastrointestinal bleeding. The most frequently implicated drug classes in these interactions include sedatives, analgesics, cardiovascular agents, and oral hypoglycaemics.

\section*{Aetiology and Pathophysiology}
The interactions between medicines and alcohol arise from distinct pharmacokinetic and pharmacodynamic mechanisms that alter how drugs are processed by the body and their biochemical effects at target sites.
\begin{itemize}
    \item \textbf{Pharmacodynamic interactions:} Additive or synergistic effects occur when alcohol and a medicine act on the same physiological system. For example, concurrent use of alcohol with other CNS depressants (such as benzodiazepines, opioids, or sedating antihistamines) leads to profound depression of the central nervous system through enhanced GABAergic transmission, potentially resulting in severe respiratory depression and death.
    \item \textbf{Acute hepatic enzyme inhibition:} Acute alcohol consumption occupies metabolic pathways, particularly CYP2E1 and other cytochrome P450 enzymes. This reduces the clearance of drugs metabolised by these pathways (such as phenytoin or warfarin), resulting in elevated plasma drug concentrations and an increased risk of toxicity.
    \item \textbf{Chronic hepatic enzyme induction:} Chronic, heavy alcohol consumption stimulates the synthesis of CYP2E1 enzymes. This accelerates the clearance of drugs metabolised by this pathway (such as paracetamol), which can reduce therapeutic efficacy when sober or generate toxic metabolites.
    \item \textbf{Paracetamol hepatotoxicity:} In the presence of chronic alcohol intake, the induced CYP2E1 pathway converts paracetamol into its highly reactive and toxic metabolite, \textit{N}-acetyl-\textit{p}-benzoquinone imine (NAPQI). Concurrently, alcohol depletes intracellular glutathione stores, which are necessary to detoxify NAPQI, leading to severe, acute centrilobular hepatic necrosis.
    \item \textbf{Vasodilation and cardiovascular effects:} Alcohol causes peripheral vasodilation and can interact with antihypertensive agents (such as beta-blockers, ACE inhibitors, and calcium channel blockers), leading to orthostatic hypotension, dizziness, and reflex tachycardia.
    \item \textbf{Gastrointestinal mucosal damage:} Alcohol directly irritates the gastric mucosa and inhibits prostaglandin synthesis. When taken with non-steroidal anti-inflammatory drugs (NSAIDs) or aspirin, this combination synergistically increases the risk of gastritis, peptic ulceration, and major gastrointestinal haemorrhage.
\end{itemize}

\section*{Clinical Features}
The clinical presentation of a drug-alcohol interaction depends on the class of medicine involved, the dose of both substances, and whether the alcohol exposure is acute or chronic.
\begin{itemize}
    \item \textbf{Central nervous system depression:} Characterised by severe drowsiness, somnolence, confusion, impaired motor coordination, slurred speech, ataxia, and in severe cases, stupor or coma.
    \item \textbf{Respiratory depression:} Depressed breathing rate, shallow respirations, hypoxia, and cyanosis, commonly observed when alcohol is combined with opioids or benzodiazepines.
    \item \textbf{Orthostatic hypotension:} Dizziness, lightheadedness, syncope, and visual disturbances when standing up, resulting from the combined vasodilatory effects of alcohol and blood pressure medications.
    \item \textbf{Hypoglycaemia:} Sweating, tremors, palpitations, confusion, and anxiety, which can occur when alcohol suppresses gluconeogenesis in patients taking insulin or oral sulfonylureas.
    \item \textbf{Gastrointestinal distress:} Epigastric pain, nausea, vomiting, haematemesis, or melaena, indicating gastric mucosal erosion or active bleeding from NSAID-alcohol interaction.
    \item \textbf{Disulfiram-like reaction:} Severe flushing, throbbing headache, tachycardia, tachypnoea, diaphoresis, nausea, and vomiting, caused by the accumulation of acetaldehyde when alcohol is consumed with medicines like metronidazole or sulfonylureas.
\end{itemize}

\section*{Differential Diagnosis}
Clinicians must distinguish drug-alcohol interactions from other acute medical emergencies, metabolic disturbances, or primary clinical presentations that present with similar symptoms.
\begin{itemize}
    \item \textbf{Accidental or intentional overdose:} Differentiating a simple interaction from a deliberate overdose of prescription medicines or illicit drugs requires a detailed history, toxidrome evaluation, and collateral information.
    \item \textbf{Acute stroke or transient ischaemic attack:} Neurological deficits, slurred speech, and ataxia caused by a cerebrovascular event must be excluded immediately through neuroimaging (such as a CT scan) and clinical examination.
    \item \textbf{Diabetic ketoacidosis or severe hypoglycaemia:} Metabolic crises in patients with diabetes can present with altered consciousness, confusion, and tachycardia, which must be differentiated from alcohol intoxication or interaction using point-of-care blood glucose and ketone testing.
    \item \textbf{Sepsis or systemic infection:} Patients presenting with confusion, tachycardia, hypotension, and tachypnoea must be screened for severe infections rather than assuming symptoms are solely due to an adverse substance interaction.
    \item \textbf{Hepatic encephalopathy:} In patients with pre-existing cirrhosis, altered mental status and asterixis must be differentiated from acute drug-alcohol toxicity through serum ammonia levels and clinical signs of liver failure.
\end{itemize}

\section*{When to Seek Medical Care}
Individuals experiencing mild symptoms of a medicine and alcohol interaction, such as mild drowsiness or dizziness, should rest, stop consuming alcohol, and contact their pharmacist or general practitioner for advice. Immediate emergency medical intervention is required if signs of severe toxicity or physiological instability develop.

\textbf{Call 000 (or local emergency services) for:}
\begin{itemize}
    \item Slow, shallow, or irregular breathing, or periods where breathing stops
    \item Complete loss of consciousness or inability to be woken up
    \item Persistent chest pain, rapid or irregular heartbeat, or severe shortness of breath
    \item Continuous vomiting, vomiting blood, or passing dark, tarry stools
    \item Sudden, severe confusion, hallucinations, or seizures
\end{itemize}

\section*{Diagnosis and Investigations}
Diagnosis of drug-alcohol interactions relies on a detailed clinical history, physical assessment, and targeted laboratory and diagnostic testing to confirm substance exposure and assess organ damage.
\begin{itemize}
    \item \textbf{Clinical history and medication review:} Obtaining a precise timeline of alcohol ingestion, along with a comprehensive list of all prescription, over-the-counter, and complementary medicines used.
    \item \textbf{Blood alcohol concentration (BAC):} Measuring serum ethanol levels to confirm ingestion and quantify the degree of intoxication.
    \item \textbf{Serum paracetamol and salicylate levels:} Checked in cases of suspected or unknown overdose to determine the risk of delayed hepatotoxicity or acid-base disturbances.
    \item \textbf{Liver and renal function tests:} Assessing serum transaminases, bilirubin, albumin, urea, and creatinine to evaluate acute liver injury or renal impairment.
    \item \textbf{Electrocardiogram (ECG):} Performed to monitor for QTc prolongation, arrhythmias, or ischemic changes, particularly if cardiotoxic medicines were co-ingested.
\end{itemize}

\section*{Management}
The management of patients with adverse medicine and alcohol interactions focuses on stabilising vital functions, preventing further absorption, managing organ-specific toxicity, and educating patients on future prevention.
\begin{itemize}
    \item \textbf{Airway protection and respiratory support:} Ensuring a patent airway, placing the patient in the recovery position to prevent aspiration of vomit, and administering oxygen or mechanical ventilation if respiratory depression is severe.
    \item \textbf{Intravenous fluid resuscitation:} Administering isotonic saline to correct hypotension, improve renal perfusion, and facilitate the excretion of metabolic waste products.
    \item \textbf{Administration of specific antidotes:} Giving naloxone for opioid-induced respiratory depression, flumazenil (with extreme caution due to seizure risks) for benzodiazepine toxicity, or acetylcysteine for paracetamol overdose.
    \item \textbf{Monitoring and management of hypoglycaemia:} Intravenous administration of dextrose (accompanied by thiamine to prevent Wernicke's encephalopathy in chronic drinkers) to correct low blood sugar.
    \item \textbf{Patient education and counselling:} Educating the patient on the importance of reading medication labels, asking pharmacists about alcohol warnings, and advising on safe drinking limits or abstinence while on therapy.
\end{itemize}

\section*{Complications}
\begin{itemize}
    \item \textbf{Respiratory arrest and hypoxic brain injury:} Caused by severe CNS and respiratory depression, which can result in permanent neurological deficit or death.
    \item \textbf{Acute liver failure:} Resulting from paracetamol toxicity or severe drug-induced liver injury, potentially requiring liver transplantation.
    \item \textbf{Severe gastrointestinal haemorrhage:} Massive bleeding from gastric ulcers or mucosal erosions, leading to haemorrhagic shock and death.
    \item \textbf{Cardiovascular collapse:} Refractory hypotension, cardiac arrhythmias, or acute myocardial infarction precipitated by hemodynamic instability.
    \item \textbf{Aspiration pneumonia:} Caused by the inhalation of gastric contents into the lungs during periods of impaired consciousness, leading to severe respiratory infection.
\end{itemize}

\section*{Prognosis and Outcomes}
The prognosis for individuals experiencing a medicine and alcohol interaction is highly variable, depending on the specific drug class involved, the dose of both substances, and the speed of medical intervention. Most mild interactions, such as transient dizziness or drowsiness, resolve completely within 12 to 24 hours as ethanol and the interacting medication are metabolised and cleared from the body. However, severe interactions involving respiratory depression, acute liver failure, or massive gastrointestinal bleeding carry a significant risk of morbidity and mortality if treatment is delayed. Long-term outcomes are generally favourable if vital organs are protected during the acute phase and if patients receive appropriate education or addiction support to prevent future occurrences. In older adults or patients with pre-existing hepatic or renal dysfunction, recovery may be prolonged, and the risk of complications like falls or permanent organ injury is higher.

\section*{Prevention and Self-Care}
\begin{itemize}
    \item \textbf{Check warning labels:} Always read prescription and over-the-counter medicine labels carefully for warnings regarding alcohol consumption.
    \item \textbf{Consult healthcare professionals:} Ask a doctor or pharmacist about potential alcohol interactions before starting any new medication, including herbal supplements.
    \item \textbf{Avoid alcohol during acute therapy:} Abstain from drinking alcohol entirely when taking short-term medicines with high interaction risks, such as antibiotics (such as metronidazole) or strong pain relievers.
    \item \textbf{Limit alcohol with chronic medicines:} If taking long-term medications for conditions like blood pressure or diabetes, discuss safe, minimal alcohol limits with your treating doctor.
    \item \textbf{Inform your care team:} Keep an updated list of all medications you take and share it with your healthcare providers, and be honest about your usual level of alcohol consumption.
\end{itemize}

\section*{Sources and Further Reading}
The following reputable organisations provide further information on medicines and alcohol safety:
\begin{itemize}
    \item Healthdirect Australia. \textit{Detailed guide on medicines and alcohol interactions.} https://www.healthdirect.gov.au/
    \item Better Health Channel. \textit{Information on how alcohol affects different types of medicines.} https://www.betterhealth.vic.gov.au/
    \item Australian Government Department of Health and Aged Care. \textit{Resources on alcohol and prescription medicine safety.} https://www.health.gov.au/
    \item NHS. \textit{Guidance on drinking alcohol while taking common prescription medicines.} https://www.nhs.uk/
    \item Mayo Clinic. \textit{Understanding the dangers of mixing alcohol with prescription and OTC drugs.} https://www.mayoclinic.org/
\end{itemize}
